El diseño de la interacción partió de una regla extraída del propio contrato
(\cref{SEC:AN-BACKEND}): las pantallas se ordenan según el ciclo de vida del recurso
principal, la instancia de examen, que recorre los estados \texttt{ASSEMBLING},
\texttt{RECEIVED}, \texttt{QUEUED}, \texttt{PENDING\_GRADING}, \texttt{GRADED},
\texttt{PUBLISHED}, \texttt{IN\_REVIEW}, \texttt{FINALIZED} y \texttt{ARCHIVED}, con el
desvío \texttt{PENDING\_REVIEW} cuando una revisión obliga a recalificar. Cada transición
permitida se convierte en una acción con su permiso y su error esperado; cada listado
paginado, en una tabla con estados de carga, vacío, error y falta de permiso; cada condición
de visibilidad del servidor, en una variante de la vista.

\subsection{Una ruta, dos vistas según el rol}

Como el rol depende de la asignatura, la ruta del examen (\texttt{/exams/:id}) no es la
misma pantalla para todos: si el usuario es estudiante en esa asignatura ve su resultado;
si es docente o coordinador, el panel del examen. Esta decisión evita duplicar rutas y hace
que un mismo enlace al examen lleve a cada persona a lo que le corresponde.

\subsection{Flujo docente}

El recorrido principal de un docente es: inicio $\rightarrow$ asignatura $\rightarrow$
examen $\rightarrow$ corrección $\rightarrow$ marcar como corregido $\rightarrow$ publicar
notas $\rightarrow$ abrir la revisión $\rightarrow$ resolver solicitudes.
\begin{itemize}
  \item \textbf{Inicio y asignaturas.} La página de inicio resume las tareas pendientes
        del usuario (exámenes por corregir y revisiones por resolver), tal como las
        devuelve la operación \texttt{/my-tasks/}. Cada asignatura muestra sus exámenes y, si el usuario
        tiene permisos de gestión, enlaza con la administración de miembros y grupos.
  \item \textbf{Panel del examen.} Lista las instancias con su estado como insignia de
        color, filtros por estado (pendientes, corregidas, publicadas) y las acciones que
        los permisos del usuario autorizan: publicar notas, gestionar la revisión, exportar
        notas, resolver incidencias de escaneo y configurar el examen. Solo una acción por
        vista se presenta como principal.
  \item \textbf{Configuración del examen.} Permite crear modelos, subir el \acs{pdf} en
        blanco, definir problemas con sus \dfnpl{rubrica}, programar la convocatoria y
        establecer las reglas de asignación de correctores.
  \item \textbf{Corrección.} La pantalla de corrección sitúa el examen escaneado en el
        centro, sobre un fondo de mesa que lo distingue de la interfaz, y a un lado los
        problemas con sus criterios de \dfn{rubrica} o su nota manual. Se navega entre
        páginas con las flechas del teclado o avance y retroceso de página, y se deshace el
        último trazo con \texttt{Ctrl+Z}. Las notas se envían con el valor anterior para el
        control de concurrencia: si otra persona ha calificado el mismo problema entretanto,
        el servidor responde \texttt{VERSION\_CONFLICT} y la interfaz lo explica en lugar de
        sobrescribir. Las anotaciones se guardan solas (\cref{SEC:DES-ANNOTATION}) y la
        barra de herramientas muestra en todo momento si el trabajo está guardado.
  \item \textbf{Incidencias de escaneo.} Cuando el reconocimiento no puede asignar una
        página o un examen a un estudiante, el panel enlaza con una pantalla que permite
        identificar al estudiante y el modelo, reordenar, aceptar, mover o descartar
        páginas y adjuntar páginas huérfanas.
  \item \textbf{Revisiones.} Una bandeja reúne las solicitudes de revisión pendientes que
        el servidor asigna al usuario; cada una se resuelve con un mensaje que el estudiante
        verá. La bandeja se oculta a quien solo es estudiante en todas sus asignaturas.
\end{itemize}

\subsection{Flujo del estudiante}

El recorrido del estudiante es más corto: inicio $\rightarrow$ examen $\rightarrow$ nota y
desglose $\rightarrow$ solicitud de revisión. La vista se diseñó para leerse de arriba
abajo: primero la nota, grande y en tipografía con serifa; después, el estado del periodo
de revisión con sus fechas; luego, si la asignatura lo permite, el desglose por problema,
las páginas corregidas con las anotaciones y el comentario general del profesorado; y por
último el formulario de revisión, en el que se marcan los problemas cuya corrección se
quiere revisar y se explica el motivo. Cada bloque depende de los permisos de visibilidad
que el examen concede al estudiante y de la ventana de revisión, tal como los aplica el
\dfn{backend}; cuando un bloque no está disponible, la vista dice por qué y cuándo lo
estará, en lugar de omitirlo sin más.

\subsection{Gestión de la organización}

Los gestores de la organización (\texttt{is\_staff}) disponen de rutas protegidas por una
guarda específica: usuarios, cursos, configuración y auditoría, implementadas en la línea
paralela (\cref{SEC:IMP-PARALLEL}), y la ingesta manual de escaneos con la consulta de
motores de reconocimiento y páginas sin examen.

\subsection{Principios comunes}

Todas las pantallas comparten unas reglas de composición: empiezan con una cabecera con
migas de pan, un antetítulo con contexto real (asignatura, fecha), el título y, si procede,
una entradilla; los registros se muestran en tablas y no en mosaicos de tarjetas idénticas;
las acciones destructivas piden confirmación en un diálogo; los estados vacíos dicen qué
va a ocurrir a continuación; y el mismo verbo acompaña a la acción desde el botón hasta el
aviso de confirmación (``Publicar notas'' $\rightarrow$ ``Notas publicadas'').

\pendiente{Incluir capturas de pantalla (files/) del panel del examen, la corrección y la
vista del estudiante, y un diagrama de flujo por actor.}
